
\newif\ifglobal

%%%%%%%%%%%%%%%%%%%%%%%
%%% Compile to view %%%
%%%%%%%%%%%%%%%%%%%%%%%

%%%%%%%%%%%%%%%%%%%%%%%%%%%%%%%%%%%%%%%%%%%%%%%%%%%
%%% Readme:                                     %%%
%%% https://www.overleaf.com/read/bwdjvfjksvjy/ %%%
%%%%%%%%%%%%%%%%%%%%%%%%%%%%%%%%%%%%%%%%%%%%%%%%%%%

% >>> M?: marks a module setting number or important notice

% >>> M4: Set {./relative-path} to external files for this module <<<
\def \modulefiles {./42-MIPIDSI-latest}
\def \modulefigures {./42-MIPIDSI/figures}
% To add an external file, use:
% (Replace '---' with actual filename)
% '\input{\modulefiles/tables/---.tex}' for tables
% '\subfile{\modulefiles/---}' for register files


% >>> M1: This module is in global project? <<<
% 'YES' - uncomment; 'NO' - comment out
\globaltrue

\ifglobal
    \documentclass[main\_\_CN.tex]{subfiles}
   
\else
    \documentclass[10pt, openany]{book}

    % >>> M2: In ./00-shared/preamble-trm-module, also find
    %         settings for visibility of todo notes, watermark <<<
    \usepackage{preamble-trm-module}


    % Enable Chinese typesetting
    \usepackage{ctex}

    % >>> M3: Temporary labels in English,
    %         you can add your own labels there <<<
    \myexternaldocument{./00-shared/temp-labels-cn}

\fi %\ifglobal


\setcounter{secnumdepth}{4}


% >>> M5: If this module needs any updates in
% './00-shared/preamble-trm-module.sty'
% (include additional package, etc.), contact your module PM
% or create the file './00-shared/changelog.tex' make the
% necessary updates yourself and document them in this file <<<

% Make text left-aligned, comment out for full justification
\raggedright


\begin{document}


% Sets language into which pre-defined bits of text are translated
\selectlanguage{chinese}

% Do not delete this block
\ifglobal
% Add the following front matter if
% this module is outside of global project
\else
    \InputIfFileExists{readme}{}{}
    {\let\clearpage\relax \listoftodos}
    {\let\clearpage\relax \listofdonetodos}
    \tableofcontents
    %\listoftables
    %\listoffigures
\fi


%%%%%%%%%%%%%%%%%%%%%%%%%
%%% TRM Module Begins %%%
%%%%%%%%%%%%%%%%%%%%%%%%%

\hypertarget{mipidsi}{}
\chapter{MIPI DSI}
\label{mod:mipidsi}



%%%%%%%%%%%%%%%%%%%%
%%% 引言 %%%
%%%%%%%%%%%%%%%%%%%%
\section{介绍}
MIPI DSI (Mobile Industry Processor Interface Display Serial Interface) 代表移动行业处理器接口显示串行接口，是一个标准接口协议，用于将显示器连接到移动设备，如智能手机、平板电脑和其他嵌入式系统。在 \chipseries{} 中，MIPI DSI 由三个部分组成：一个 DSI 桥接器，一个 MIPI DSI 主控制器，以及一个 MIPI 双向 D-PHY。它提供一个时钟通道和两条数据通道，并支持每个通道最高 1.5 Gbps 的通信速率，支持与符合 DSI 标准的显示器通信。DSI 属于 MIPI 联盟定义的一组通信协议。

DSI 桥接器提供系统内存与 MIPI DSI 主控制器之间的接口。它通过 VDMA 从系统内存接收图像数据，并通过显示像素接口 (DPI) 将图像数据输出到 MIPI DSI 主控制器中。
MIPI DSI 主控制器实现了 MIPI DSI 规范中定义的协议功能。它从 DSI 桥接器接收图像数据，并通过 PHY 协议接口 (PPI) 将图像数据输出到 MIPI 双向 D-PHY。
MIPI 双向 D-PHY 实现了 MIPI D-PHY 接口的双向通用通道的物理层，提供一个时钟通道和两个数据通道。PHY 仅用作主设备。

%%%%%%%%%%%%%%%%%%%%
%%% Terminology %%%
%%%%%%%%%%%%%%%%%%%%



%%%%%%%%%%%%%%%%%%%%
%%% Feature List %%%
%%%%%%%%%%%%%%%%%%%%

\section{功能列表}
\label{sec:mipidsi-features}

MIPI DSI 在 \chipname{} 中支持以下功能：
\begin{itemize}
    \item 在从内存读取图像数据时，将 DSI 桥或 VDMA 配置为流控制器
    \item 支持在 DSI 桥中配置 DPI 视频时序
    \item 支持多种输入图像格式：RGB888、RGB666、RGB565、YUV422、YUV420、GRAY
    \item 支持各种 DPI 图像输出格式：RGB888、RGB666、RGB565
    \item 最大 DPI 频率：240 MHz
    \item MIPI DSI 主机符合 MIPI DSI 规范
    \item 符合 MIPI D-PHY 接口规范，修订版 1.1 的 PHY
    \item PHY 包含 PLL
    \item 连续和非连续时钟通道模式
    \item 支持通过数据通道 0 进入逃逸模式
    \item 在 HS TX 中，通道操作范围为 80 Mbps 到 1.5 Gbps
    \item 支持突发和非突发视频模式
    \item 支持通过 APB 接口的命令模式访问显示寄存器
    \item 支持 MIPI 联盟规范中定义的显示命令集 (DCS)
    \item 视频图案生成器
    \item 支持超低功耗模式
\end{itemize}


%%%%%%%%%%%%%%%%%%%%%%%%%%%%%%
%%% 架构概述 %%%
%%%%%%%%%%%%%%%%%%%%%%%%%%%%%%

\section{架构概述}
\label{sec:mipidsi-arch-overview}
图 \ref{figure:mipidsi_arch} 显示了 MIPI DSI 的架构。

\begin{figure}[H]
    \centering
    \includegraphics[width=1.0\textwidth]{\modulefigures/mipidsi_arch_CN.pdf}
    \caption{MIPI DSI 架构}
    \label{figure:mipidsi_arch}
\end{figure}
\begin{itemize}
    \item \textbf{DSI 桥接器}：DSI 桥接器通过 VDMA 从系统内存接收图像数据，将其转换为 DPI 数据格式，然后输出到 MIPI DSI 主机。
    \item \textbf{MIPI DSI 主机}：MIPI DSI 主机从两个来源接收图像数据：DSI 桥接器或视频图案生成器。然后，接收到的图像数据由数据包处理器根据 MIPI DSI 规范进行打包，并发送到 PHY 接口控制器。最后，图像数据被转换为 PPI 数据格式并输出到 MIPI DSI 双向 D-PHY。MIPI DSI 主机包括：

    \begin{itemize}
        \item \textbf{DPI 接口}：从 DSI 桥接器接收像素数据并将其缓冲到 FIFO。
        \item \textbf{寄存器组}：提供对配置和控制寄存器的访问。
        \item \textbf{APB 到通用接口}：将 APB 操作转换为通用命令。
        \item \textbf{视频图案生成器}：生成视频测试图案。
        \item \textbf{数据包处理器}：安排链路内部的活动，包括数据包构建、生成 ECC CRC 代码、数据包接收等。
        \item \textbf{PHY 接口控制器}：管理 PHY PPI 接口，包括启用低功耗 (LP) 传输/接收、高速 (HS) 传输等。
        \item \textbf{错误管理}：监控并通知 DSI 链路上的错误状态。
    \end{itemize}
    \item \textbf{MIPI DSI D-PHY}：MIPI DSI 双向 D-PHY 实现了具有一个时钟通道和两个数据通道的物理层，连接到外部显示驱动器。MIPI DSI 双向 D-PHY 还包括一个 PLL 以生成 HS TX 时钟。
\end{itemize}


%%%%%%%%%%%%%%%
%%% 功能描述 %%%
%%%%%%%%%%%%%%%

\section{功能描述}
\label{sec:mipidsi-func-descr}
本节介绍了 MIPI 双向 (BD) D-PHY 的不同操作状态和模式，以及一些有助于 MIPI DSI 使用的主要功能。

\subsection{DSI 桥接器}
\label{sec:dsi-bridge}

DSI Bridge 通过 VDMA 从系统内存接收图像数据。数据传输流控制器可以通过 \hyperref[fielddesc:DSIBRGDSIDMAFLOWCONTROLLER]{DSI\_BRG\_DSI\_DMA\_FLOW\_CONTROLLER} 配置为 VDMA 或 DSI Bridge。建议使用 VDMA 作为流控制器以获得更好的性能。




当使用 VDMA 作为流控制器接收图像数据时，可以将图像拆分为多个块传输。要启用此功能，请将 \hyperref[fielddesc:DSIBRGDMAMULTIBLKEN]{DSI\_BRG\_DMA\_MULTIBLK\_EN} 设置为 1，并将 \hyperref[fielddesc:DSIBRGDMAFLOWMULTIBLKNUM]{DSI\_BRG\_DMA\_FLOW\_MULTIBLK\_NUM} 配置为所需值。有关块传输的详细说明，请参见 \ref{mod:vdma} \textit{\nameref{mod:vdma}} 的 \textit{传输层次结构} 小节。

DSI Bridge 支持 RGB888、RGB666、RGB565、YUV422、YUV420 和 GRAY 数据类型的输入。对于 YUV422，DSI Bridge 支持 8 种不同的像素排列。用户可以通过配置 \hyperref[fielddesc:DSIBRGDATAINTYPE]{DSI\_BRG\_DATA\_IN\_TYPE} 寄存器来选择输入数据类型是否为 YUV，然后使用 \hyperref[fielddesc:DSIBRGRAWTYPE]{DSI\_BRG\_RAW\_TYPE} 指定具体的输入格式。对于发送到 DSI Host 的输出格式，使用 \hyperref[fielddesc:DSIBRGDPITYPE]{DSI\_BRG\_DPI\_TYPE} 指定 RGB 格式。

所有支持的输入像素排列，参见表 \ref{tab:input-pixel-arrangement}。


\scriptsize
\begin{longtable}[c]{ |p{2cm}|p{1.5cm}|p{1.5cm}|p{1.5cm}|p{1.5cm}|p{1.5cm}|p{1.5cm}|p{1.5cm}|p{1.5cm}| }
    \caption{输入像素排列}
    \label{tab:input-pixel-arrangement}\\\hline

    \rowcolor{lightgray}
    \textbf{颜色模式} & \textbf{Addr + 7} & \textbf{Addr + 6} & \textbf{Addr + 5} & \textbf{Addr + 4} & \textbf{Addr + 3} & \textbf{Addr + 2} & \textbf{Addr + 1} & \textbf{Addr + 0} \\\hline
    \endfirsthead


    \multicolumn{9}{c}{\bfseries \tablename\ \thetable{} -- 接上页}\\\hline
    \rowcolor{lightgray}
    \textbf{颜色模式} & \textbf{Addr + 7} & \textbf{Addr + 6} & \textbf{Addr + 5} & \textbf{Addr + 4} & \textbf{Addr + 3} & \textbf{Addr + 2} & \textbf{Addr + 1} & \textbf{Addr + 0} \\\hline
    \endhead

    \multicolumn{9}{r}{\textbf{见下页}} \\
    \endfoot

    \endlastfoot
        \multirow{3}{*}{\textbf{RGB888}} & \textbf{G2}[7:0] & \textbf{B2}[7:0] & \textbf{R1}[7:0] & \textbf{G1}[7:0] & \textbf{B1}[7:0] & \textbf{R0}[7:0] & \textbf{G0}[7:0] & \textbf{B0}[7:0] \\\cline{2-9}
        & \textbf{B5}[7:0] & \textbf{R4}[7:0] & \textbf{G4}[7:0] & \textbf{B4}[7:0] & \textbf{R3}[7:0] & \textbf{G3}[7:0] & \textbf{B3}[7:0] & \textbf{R2}[7:0] \\\cline{2-9}
        & \textbf{R7}[7:0] & \textbf{G7}[7:0] & \textbf{B7}[7:0] & \textbf{R6}[7:0] & \textbf{G6}[7:0] & \textbf{B6}[7:0] & \textbf{R5}[7:0] & \textbf{G5}[7:0] \\\hline
        \textbf{RGB565} & \textbf{R3}[4:0]\textbf{G3}[5:3] & \textbf{G3}[2:0]\textbf{B3}[4:0] & \textbf{R2}[4:0]\textbf{G2}[5:3] & \textbf{G2}[2:0]\textbf{B2}[4:0]  & \textbf{R1}[4:0]\textbf{G1}[5:3] & \textbf{G1}[2:0]\textbf{B1}[4:0] & \textbf{R0}[4:0]\textbf{G0}[5:3] & \textbf{G0}[2:0]\textbf{B0}[4:0] \\\hline
        \makecell[l]{\textbf{YUV422}\\ \hyperref[fielddesc:DSIBRGYUV422FORMAT]{\makecell[l]{DSI\_BRG\_YUV422\\\_FORMAT = 0}}\\ \hyperref[fielddesc:DSIBRGYUVPIXENDIAN]{\makecell[l]{DSI\_BRG\_YUV\\\_PIX\_ENDIAN = 0}} } & \textbf{Y3}[7:0] & \textbf{U3}[7:0] & \textbf{Y2}[7:0] & \textbf{V2}[7:0] & \textbf{Y1}[7:0] & \textbf{U1}[7:0] & \textbf{Y0}[7:0] & \textbf{V0}[7:0] \\\hline
        \makecell[l]{\textbf{YUV422}\\ \hyperref[fielddesc:DSIBRGYUV422FORMAT]{\makecell[l]{DSI\_BRG\_YUV422\\\_FORMAT = 0}}\\ \hyperref[fielddesc:DSIBRGYUVPIXENDIAN]{\makecell[l]{DSI\_BRG\_YUV\\\_PIX\_ENDIAN = 1}}  } & \textbf{Y0}[7:0] & \textbf{U0}[7:0] & \textbf{Y1}[7:0] & \textbf{V1}[7:0] & \textbf{Y2}[7:0] & \textbf{U2}[7:0] & \textbf{Y3}[7:0] & \textbf{V3}[7:0] \\\hline
        \makecell[l]{\textbf{YUV422}\\ \hyperref[fielddesc:DSIBRGYUV422FORMAT]{\makecell[l]{DSI\_BRG\_YUV422\\\_FORMAT = 1}} \\ \hyperref[fielddesc:DSIBRGYUVPIXENDIAN]{\makecell[l]{DSI\_BRG\_YUV\\\_PIX\_ENDIAN = 0}}  } & \textbf{Y3}[7:0] & \textbf{V3}[7:0] & \textbf{Y2}[7:0] & \textbf{U2}[7:0] & \textbf{Y1}[7:0] & \textbf{V1}[7:0] & \textbf{Y0}[7:0] & \textbf{U0}[7:0] \\\hline
        \makecell[l]{\textbf{YUV422}\\ \hyperref[fielddesc:DSIBRGYUV422FORMAT]{\makecell[l]{DSI\_BRG\_YUV422\\\_FORMAT = 1}} \\ \hyperref[fielddesc:DSIBRGYUVPIXENDIAN]{\makecell[l]{DSI\_BRG\_YUV\\\_PIX\_ENDIAN = 1}}  } & \textbf{Y0}[7:0] & \textbf{V0}[7:0] & \textbf{Y1}[7:0] & \textbf{U1}[7:0] & \textbf{Y2}[7:0] & \textbf{V2}[7:0] & \textbf{Y3}[7:0] & \textbf{U3}[7:0] \\\hline
        \makecell[l]{\textbf{YUV422}\\ \hyperref[fielddesc:DSIBRGYUV422FORMAT]{\makecell[l]{DSI\_BRG\_YUV422\\\_FORMAT = 2}} \\ \hyperref[fielddesc:DSIBRGYUVPIXENDIAN]{\makecell[l]{DSI\_BRG\_YUV\\\_PIX\_ENDIAN = 0}}  } & \textbf{U3}[7:0] & \textbf{Y3}[7:0] & \textbf{V2}[7:0] & \textbf{Y2}[7:0] & \textbf{U1}[7:0] & \textbf{Y1}[7:0] & \textbf{V0}[7:0] & \textbf{Y0}[7:0] \\\hline
        \makecell[l]{\textbf{YUV422}\\ \hyperref[fielddesc:DSIBRGYUV422FORMAT]{\makecell[l]{DSI\_BRG\_YUV422\\\_FORMAT = 2}} \\ \hyperref[fielddesc:DSIBRGYUVPIXENDIAN]{\makecell[l]{DSI\_BRG\_YUV\\\_PIX\_ENDIAN = 1}}  } & \textbf{U0}[7:0] & \textbf{Y0}[7:0] & \textbf{V1}[7:0] & \textbf{Y1}[7:0] & \textbf{U2}[7:0] & \textbf{Y2}[7:0] & \textbf{V3}[7:0] & \textbf{Y3}[7:0] \\\hline
        \makecell[l]{\textbf{YUV422}\\ \hyperref[fielddesc:DSIBRGYUV422FORMAT]{\makecell[l]{DSI\_BRG\_YUV422\\\_FORMAT = 3}} \\ \hyperref[fielddesc:DSIBRGYUVPIXENDIAN]{\makecell[l]{DSI\_BRG\_YUV\\\_PIX\_ENDIAN = 0}} } & \textbf{V3}[7:0] & \textbf{Y3}[7:0] & \textbf{U2}[7:0] & \textbf{Y2}[7:0] & \textbf{V1}[7:0] & \textbf{Y1}[7:0] & \textbf{U0}[7:0] & \textbf{Y0}[7:0] \\\hline
        \makecell[l]{\textbf{YUV422}\\ \hyperref[fielddesc:DSIBRGYUV422FORMAT]{\makecell[l]{DSI\_BRG\_YUV422\\\_FORMAT = 3}} \\ \hyperref[fielddesc:DSIBRGYUVPIXENDIAN]{\makecell[l]{DSI\_BRG\_YUV\\\_PIX\_ENDIAN = 1}}  } & \textbf{V0}[7:0] & \textbf{Y0}[7:0] & \textbf{U1}[7:0] & \textbf{Y1}[7:0] & \textbf{V2}[7:0] & \textbf{Y2}[7:0] & \textbf{U3}[7:0] & \textbf{Y3}[7:0] \\\hline
        \multirow{3}{*}{\makecell[l]{\textbf{YUV420}\\ Odd lines (1, 3, 5...)}} & \textbf{Y4}[7:0] & \textbf{U4}[7:0] & \textbf{Y3}[7:0] & \textbf{Y2}[7:0] & \textbf{U2}[7:0] & \textbf{Y1}[7:0] & \textbf{Y0}[7:0] & \textbf{U0}[7:0] \\\cline{2-9}
        & \textbf{U10}[7:0] & \textbf{Y9}[7:0] & \textbf{Y8}[7:0] & \textbf{U8}[7:0] & \textbf{Y7}[7:0] & \textbf{Y6}[7:0] & \textbf{U6}[7:0] & \textbf{Y5}[7:0] \\\cline{2-9}
        & \textbf{Y15}[7:0] & \textbf{Y14}[7:0] & \textbf{U14}[7:0] & \textbf{Y13}[7:0] & \textbf{Y12}[7:0] & \textbf{U12}[7:0] & \textbf{Y11}[7:0] & \textbf{Y10}[7:0] \\\hline

        \multirow{3}{*}{\makecell[l]{\textbf{YUV420}\\ Even lines (2, 4, 6...)}} & \textbf{Y4}[7:0] & \textbf{V4}[7:0] & \textbf{Y3}[7:0] & \textbf{Y2}[7:0] & \textbf{V2}[7:0] & \textbf{Y1}[7:0] & \textbf{Y0}[7:0] & \textbf{V0}[7:0] \\\cline{2-9}
        & \textbf{V10}[7:0] & \textbf{Y9}[7:0] & \textbf{Y8}[7:0] & \textbf{V8}[7:0] & \textbf{Y7}[7:0] & \textbf{Y6}[7:0] & \textbf{V6}[7:0] & \textbf{Y5}[7:0] \\\cline{2-9}
        & \textbf{Y15}[7:0] & \textbf{Y14}[7:0] & \textbf{V14}[7:0] & \textbf{Y13}[7:0] & \textbf{Y12}[7:0] & \textbf{V12}[7:0] & \textbf{Y11}[7:0] & \textbf{Y10}[7:0] \\\hline

        \makecell[l]{\textbf{GRAY}} & \textbf{P7}[7:0] & \textbf{P6}[7:0] & \textbf{P5}[7:0] & \textbf{P4}[7:0] & \textbf{P3}[7:0] & \textbf{P2}[7:0] & \textbf{P1}[7:0] & \textbf{P0}[7:0] \\\hline


\end{longtable}
\normalsize

DSI Bridge 通过 DPI 接口输出数据，并支持由 \hyperref[fielddesc:DSIBRGDPICONFIG]{DSI\_BRG\_DPI\_CONFIG} 指定的几种 RGB 输出位排列。所有支持的位排列，请参见表 \ref{tab:dpi-bit-arrangement}。

\scriptsize
\begin{longtable}[c]{ |p{2.5cm}|p{1cm}|p{1cm}|p{1cm}|p{1cm}|p{1cm}|p{1cm}| }
    \caption{DPI 位排列}
    \label{tab:dpi-bit-arrangement}\\\hline

    \rowcolor{lightgray}
    \textbf{} & \multicolumn{3}{|c|}{\textbf{RGB565}} & \multicolumn{2}{|c|}{\textbf{RGB666}} & \textbf{RGB888}  \\\hline
    \rowcolor{lightgray}
    \textbf{\hyperref[fielddesc:DSIBRGDPICONFIG]{DSI\_BRG\_DPI\_CONFIG}} & \textbf{0} & \textbf{1} & \textbf{2} & \textbf{0} & \textbf{1} & \textbf{0}  \\\hline
    \endfirsthead


    \multicolumn{7}{c}{\bfseries \tablename\ \thetable{} -- 接上页}\\\hline
    \rowcolor{lightgray}
    \textbf{} & \multicolumn{3}{|c|}{\textbf{RGB565}} & \multicolumn{2}{|c|}{\textbf{RGB666}} & \textbf{RGB888}  \\\hline
    \rowcolor{lightgray}
    \textbf{\hyperref[fielddesc:DSIBRGDPICONFIG]{DSI\_BRG\_DPI\_CONFIG}} & \textbf{0} & \textbf{1} & \textbf{2} & \textbf{0} & \textbf{1} & \textbf{0} \\\hline
    \endhead

    \multicolumn{7}{r}{\textbf{见下页}} \\
    \endfoot

    \endlastfoot
        
        bit-23 &    &    &    &    &    & R7 \\\hline
        bit-22 &    &    &    &    &    & R6 \\\hline
        bit-21 &    &    & R4 &    & R5 & R5 \\\hline
        bit-20 &    & R4 & R3 &    & R4 & R4 \\\hline
        bit-19 &    & R3 & R2 &    & R3 & R3 \\\hline
        bit-18 &    & R2 & R1 &    & R2 & R2 \\\hline
        bit-17 &    & R1 & R0 & R5 & R1 & R1 \\\hline
        bit-16 &    & R0 &    & R4 & R0 & R0 \\\hline
        bit-15 & R4 &    &    & R3 &    & G7 \\\hline
        bit-14 & R3 &    &    & R2 &    & G6 \\\hline
        bit-13 & R2 & G5 & G5 & R1 & G5 & G5 \\\hline
        bit-12 & R1 & G4 & G4 & R0 & G4 & G4 \\\hline
        bit-11 & R0 & G3 & G3 & G5 & G3 & G3 \\\hline
        bit-10 & G5 & G2 & G2 & G4 & G2 & G2 \\\hline
        bit-9  & G4 & G1 & G1 & G3 & G1 & G1 \\\hline
        bit-8  & G3 & G0 & G0 & G2 & G0 & G0 \\\hline
        bit-7  & G2 &    &    & G1 &    & B7 \\\hline
        bit-6  & G1 &    &    & G0 &    & B6 \\\hline
        bit-5  & G0 &    & B4 & B5 & B5 & B5 \\\hline
        bit-4  & B4 & B4 & B3 & B4 & B4 & B4 \\\hline
        bit-3  & B3 & B3 & B2 & B3 & B3 & B3 \\\hline
        bit-2  & B2 & B2 & B1 & B2 & B2 & B2 \\\hline
        bit-1  & B1 & B1 & B0 & B1 & B1 & B1 \\\hline
        bit-0  & B0 & B0 &    & B0 & B0 & B0 \\\hline
        

\end{longtable}
\normalsize

\subsection{DSI 主机}

\subsubsection{视频传输}
DSI Host 使用 DPI 作为输入接口。它用于在视频模式下从 HOST 向外设传输实时像素流。

有两种视频传输模式：

\begin{itemize}
    \item \textbf{突发模式}：在突发模式下，一行像素会被缓冲到FIFO中，并在一个单独的数据包中无中断地传输。这使得像素行能够在短时间内发送出去，然后让通道回到低功耗 (LP) 状态，从而降低功耗，并为其他数据传输节省时间。
    \item \textbf{非突发模式}：在非突发模式下，一行像素被分成几个数据包。这种模式用于匹配像素带宽与 DSI 链路带宽，并减少主机和外设的内存需求。非突发模式有两种：
    \begin{itemize}
        \item 带同步脉冲的非突发模式：允许外设准确重建原始视频时序，包括同步脉冲宽度。
        \item 带同步事件的非突发模式：不需要准确重建同步脉冲。
    \end{itemize}
\end{itemize}

如果主机和外设满足以下要求，建议选择突发模式。否则，像素数据可能会丢失，导致显示外设发生故障。
\begin{itemize}
    \item 显示外设支持在一个单一块中接收一行完整像素，并且接收缓冲区不会溢出。
    \item DSI 输出带宽远高于 DPI 输入带宽，该带宽差异使传输通道能够在每一行数据传输完成后进入低功耗 (LP) 状态。
\end{itemize}

如果无法满足任何要求，则应选择非突发模式。在非突发模式下，DSI 主机应配置为匹配 DSI 输出比率和 DPI 输入比率。以下方程显示了传输参数应如何配置（左侧表示 DSI 输出比率，右侧表示 DPI 输入比率）：
\begin{itemize}
    \item 如果启用了空数据包 (\hyperref[fielddesc:DSIHOSTVIDNULLSIZE]{DSI\_HOST\_VID\_NULL\_SIZE} $\neq$ 0)：
    \begin{align*}
    & lanebyteclk\_period \times \hyperref[fielddesc:DSIHOSTVIDNUMCHUNKS]{DSI\_HOST\_VID\_NUM\_CHUNKS} \\
    & \quad \times (\hyperref[fielddesc:DSIHOSTVIDPKTSIZE]{DSI\_HOST\_VID\_PKT\_SIZE} \times bytes\_per\_pixel + 12 \\
    & \quad + \hyperref[fielddesc:DSIHOSTVIDNULLSIZE]{DSI\_HOST\_VID\_NULL\_SIZE}) / number\_of\_lanes \\
    &= dpiclk\_period \times pixels\_per\_line
    \end{align*}
    \item 如果禁用了空数据包 (\hyperref[fielddesc:DSIHOSTVIDNULLSIZE]{DSI\_HOST\_VID\_NULL\_SIZE} $=$ 0)：
    \begin{align*}
    & lanebyteclk\_period \times \hyperref[fielddesc:DSIHOSTVIDNUMCHUNKS]{DSI\_HOST\_VID\_NUM\_CHUNKS} \\
    & \quad \times (\hyperref[fielddesc:DSIHOSTVIDPKTSIZE]{DSI\_HOST\_VID\_PKT\_SIZE} \times bytes\_per\_pixel + 6) \\
    & \quad / number\_of\_lanes = dpiclk\_period \times pixels\_per\_line
    \end{align*}
\end{itemize}

有关 lanebyteclk 和 dpiclk 的定义，请参阅 \ref{sec:clocks-and-display-resolution}。

\subsubsection{命令传输}
\label{sec:command-transmission}
DSI 主机允许通过 APB 从接口，以命令模式传输 DSI 规范中厂商自定义的通用命令。

命令数据包使用 APB 寄存器构建：
\begin{itemize}
    \item \textbf{\hyperref[regdesc:DSIHOSTGENPLDDATAREG]{DSI\_HOST\_GEN\_PLD\_DATA\_REG}}：向该寄存器写入数据会将数据作为命令数据包的有效载荷发送。读取该寄存器会返回回读数据包的有效载荷。
    \item \textbf{\hyperref[regdesc:DSIHOSTGENHDRREG]{DSI\_HOST\_GEN\_HDR\_REG}}：包含命令数据包头类型和头数据。写入此寄存器会触发命令传输。
\end{itemize}

有效数据包如下：
\begin{itemize}
    \item 通用写短包 0 参数
    \item 通用写短包 1 参数
    \item 通用写短包 2 参数
    \item 通用读短包 0 参数
    \item 通用读短包 1 参数
    \item 通用读短包 2 参数
    \item 最大读包配置
    \item 通用长写包
\end{itemize}

在视频模式的消隐期 (Blanking) 或低功耗期 (BLLP) 内，命令数据包可以在高速 (HS) 或低速 (LP) 通道状态下传输。BLLP 对应于图 \ref{figure:command-transmission-periods} 中的阴影区域。

\begin{figure}[H]
    \centering
    \includegraphics[width=0.8\textwidth]{\modulefigures/command-transmission-periods.pdf}
    \caption{命令传输周期}
    \label{figure:command-transmission-periods}
\end{figure}

\paragraph{LP 模式下的命令传输}

要在视频传输期间以 LP 模式发送命令，请将 \hyperref[fielddesc:DSIHOSTLPCMDEN]{DSI\_HOST\_LP\_CMD\_EN} 设置为 1。在这种情况下，即使 \hyperref[regdesc:DSIHOSTVIDMODECFGREG]{DSI\_HOST\_VID\_MODE\_CFG\_REG} 的第 8 到 13 位设置为 0，仍然可以发送命令。

\begin{itemize}
    \item \hyperref[fielddesc:DSIHOSTOUTVACTLPCMDTIME]{DSI\_HOST\_OUTVACT\_LPCMD\_TIME} 配置在 VSA、VBP 和 VFP 期间以 LP 模式传输命令的可用时间。此时间基于逃逸时钟确定，单位为字节。非突发带同步脉冲的 outvact\_lpcmd\_time 如图 \ref{figure:outvact-lpcmd-time-for-non-burst-with-sync-pulse} 所示；突发或非突发带同步事件的 outvact\_lpcmd\_time 如图 \ref{figure:outvact-lpcmd-time-for-burst-or-non-burst-with-sync-event} 所示。
    \begin{figure}[H]
        \centering
        \includegraphics[width=1.0\textwidth]{\modulefigures/outvact-lpcmd-time-for-non-burst-with-sync-pulse.pdf}
        \caption{非突发带同步脉冲的 outvact\_lpcmd\_time}
        \label{figure:outvact-lpcmd-time-for-non-burst-with-sync-pulse}
    \end{figure}
    \begin{figure}[H]
        \centering
        \includegraphics[width=1.0\textwidth]{\modulefigures/outvact-lpcmd-time-for-burst-or-non-burst-with-sync-event.pdf}
        \caption{突发或非突发带同步事件的 outvact\_lpcmd\_time}
        \label{figure:outvact-lpcmd-time-for-burst-or-non-burst-with-sync-event}
    \end{figure}
    outvact\_lpcmd\_time 的计算公式为：
    \begin{align*}
    & outvact\_lpcmd\_time = t_L - (t_{H1} + t_{HS\rightarrow LP} + t_{LP\rightarrow HS} \\
    & \quad + t_{LPDT} + 2 \times t_{ESCCLK}) \\
    & \quad / (2 \times 8 \times t_{ESCCLK})
    \end{align*}

    \item \hyperref[fielddesc:DSIHOSTINVACTLPCMDTIME]{DSI\_HOST\_INVACT\_LPCMD\_TIME} 配置在 HFP 期间以 LP 模式传输命令的可用时间。此时间基于逃逸时钟确定，单位为字节。非突发带同步脉冲的 invact\_lpcmd\_time 如图 \ref{figure:invact-lpcmd-time-for-non-burst-with-sync-pulse} 所示；非突发带同步事件的 invact\_lpcmd\_time 如图 \ref{figure:invact-lpcmd-time-for-non-burst-with-sync-event} 所示。
    \begin{figure}[H]
        \centering
        \includegraphics[width=1.0\textwidth]{\modulefigures/invact-lpcmd-time-for-non-burst-with-sync-pulse.pdf}
        \caption{非突发带同步脉冲的 invact\_lpcmd\_time}
        \label{figure:invact-lpcmd-time-for-non-burst-with-sync-pulse}
    \end{figure}
    \begin{figure}[H]
        \centering
        \includegraphics[width=1.0\textwidth]{\modulefigures/invact-lpcmd-time-for-non-burst-with-sync-event.pdf}
        \caption{突发或非突发带同步事件的 invact\_lpcmd\_time}
        \label{figure:invact-lpcmd-time-for-non-burst-with-sync-event}
    \end{figure}

    invact\_lpcmd\_time 的计算公式为：
    \begin{align*}
    & invact\_lpcmd\_time = t_L - (t_{HSA} + t_{HBP} + t_{HACT} + t_{HS\rightarrow LP} \\
    & \quad + t_{LP\rightarrow HS} + t_{LPDT} + 2 \times t_{ESCCLK}) \\
    & \quad / (2 \times 8 \times t_{ESCCLK})
    \end{align*}
\end{itemize}


\paragraph{HS 模式下的命令传输}
如果 \hyperref[fielddesc:DSIHOSTLPCMDEN]{DSI\_HOST\_LP\_CMD\_EN} 设置为 0，则在视频传输期间以 HS 模式发送命令。在这种情况下，不需要配置任何寄存器。

\subsubsection{视频图案生成器}
\label{sec:video-pattern-generator}
视频图案生成器可以输出水平或垂直的色带或 BER 测试图案。有关图案分辨率，请参阅 \hyperref[sec:clocks-and-display-resolution]{时钟和显示分辨率}。以下寄存器字段用于配置视频图案生成器：

\begin{itemize}
    \item \hyperref[fielddesc:DSIHOSTVPGORIENTATION]{DSI\_HOST\_VPG\_ORIENTATION}：指示色带的方向
    \item \hyperref[fielddesc:DSIHOSTVPGMODE]{DSI\_HOST\_VPG\_MODE}：选择图案
    \item \hyperref[fielddesc:DSIHOSTVPGEN]{DSI\_HOST\_VPG\_EN}：启用视频图案生成器
\end{itemize}

\paragraph{色带图案}
色带有 8 个颜色：白色、黄色、青色、绿色、品红色、红色、蓝色和黑色。每个颜色的宽度可以通过将图案大小除以 8 来计算。如果大小不能被 8 整除，则最后一个颜色黑色包含剩余的大小。图 \ref{figure:vertical-color-bar} 和 \ref{figure:horizontal-color-bar} 分别显示了垂直和水平色带。

\begin{figure}[H]
    \centering
    
\includegraphics[width=0.6\textwidth]{\modulefigures/vertical-color-bar.pdf}
    \caption{垂直色带}
    \label{figure:vertical-color-bar}
\end{figure}

\begin{figure}[H]
    \centering
    \includegraphics[width=0.6\textwidth]{\modulefigures/horizontal-color-bar.pdf}
    \caption{水平色带}
    \label{figure:horizontal-color-bar}
\end{figure}

表 \ref{tab:color-bar-coding} 显示了色带图案的RGB编码。
\begin{longtable}[c]{ |p{0.5cm}|p{1.5cm}|p{1.5cm}|p{1.5cm}|p{1.5cm}|p{1.5cm}|p{1.5cm}|p{1.5cm}|p{1.5cm}| }
    \caption{色带编码}
    \label{tab:color-bar-coding}\\\hline

    \rowcolor{lightgray}
     & \textbf{白色} & \textbf{黄色} & \textbf{青色} & \textbf{绿色} & \textbf{品红色} & \textbf{红色} & \textbf{蓝色} & \textbf{黑色}  \\\hline
    \endfirsthead


    \multicolumn{9}{c}{\bfseries \tablename\ \thetable{} -- 接上页}\\\hline
    \rowcolor{lightgray}
    & \textbf{白色} & \textbf{黄色} & \textbf{青色} & \textbf{绿色} & \textbf{品红色} & \textbf{红色} & \textbf{蓝色} & \textbf{黑色}  \\\hline
    \endhead

    \multicolumn{9}{r}{\textbf{见下页}} \\
    \endfoot

    \endlastfoot
        
        \textbf{R} & 高 & 高 & 低  & 低  & 高 & 高 & 低  & 低 \\\hline
        \textbf{G} & 高 & 高 & 高 & 高 & 低  & 低  & 低  & 低 \\\hline
        \textbf{B} & 高 & 低  & 高 & 低  & 高 & 低  & 高 & 低 \\\hline

\end{longtable}

\paragraph{BER 测试图案}
BER测试图案用于测试设备 RX PHY 是否正确接收数据。生成的数据图案包括：
\begin{itemize}
    \item X 字节的 0xAA（高频图案）
    \item X 字节的 0x33（中频图案）
    \item X 字节的 0xF0（低频图案）
    \item X 字节的 0x7F（孤立 0 图案）
    \item X 字节的 0x55（高频图案）
    \item X 字节的 0xCC（中频图案）
    \item X 字节的 0x0F（低频图案）
    \item Y 字节的 0x80（孤立 1 图案）
\end{itemize}

X 和 Y 取决于图案分辨率，在大多数情况下 Y 等于 X。使用 RGB888 颜色编码和八的倍数的水平分辨率，图 \ref{figure:rgb888-ber-testing-pattern} 显示了在 DSI 显示器上出现的图案：

\begin{figure}[H]
    \centering
    \includegraphics[width=0.6\textwidth]{\modulefigures/rgb888-ber-testing-pattern.pdf}
    \caption{RGB888 BER 测试图案}
    \label{figure:rgb888-ber-testing-pattern}
\end{figure}

\subsubsection{错误报告}
\hyperref[regdesc:DSIHOSTINTST0REG]{DSI\_HOST\_INT\_ST\textcolor{red}{0/1}\_REG} 寄存器指示 DSI 主机错误报告。这些寄存器还可以触发中断以通知系统。要触发错误中断，应该清除相应的掩码寄存器 \hyperref[regdesc:DSIHOSTINTMSK0REG]{DSI\_HOST\_INT\_MSK\textcolor{red}{0/1}\_REG} 中的位。 \hyperref[regdesc:DSIHOSTINTST0REG]{DSI\_HOST\_INT\_ST\textcolor{red}{0/1}\_REG} 在读取操作后将被清除。

% 



\subsection{D-PHY}
本节描述了 D-PHY 操作模式和 PLL 配置。
\subsubsection{D-PHY 操作模式}
图 \ref{figure:phy-operating-modes} 显示了 PHY 初始化模式和运行模式。
\begin{figure}[H]
    \centering
    \includegraphics[width=0.8\textwidth]{\modulefigures/phy-operating-modes_CN.pdf}
    \caption{PHY 运行模式}
    \label{figure:phy-operating-modes}
\end{figure}

\paragraph{无电源状态 (No-Power)}
无电源 (No-Power) 状态的特征是没有对 PHY 提供任何供电电压。

对于 \chipname{} 的 MIPI DSI，当电源 VDD\_HP\_\regindex{n} (n: 0-3) 和 VDD\_MIPI\_DPHY 关闭时，MIPI BD D-PHY 处于无电源状态。

\paragraph{关闭状态 (Shutdown)}
在关闭状态下，会禁用所有模拟模块，重置所有数字逻辑，DP/DN 和 CP/CN 通道处于高阻抗 (Hi-Z)。要启用 PHY 的关闭模式，将 \hyperref[fielddesc:DSIHOSTPHYRSTZ]{DSI\_HOST\_PHY\_RSTZ} 和 \hyperref[fielddesc:DSIHOSTPHYSHUTDOWNZ]{DSI\_HOST\_PHY\_SHUTDOWNZ} 设置为 0。

在离开关闭状态之前应配置一些参数，包括要使用的通道数量、通道操作频率范围，和 PLL 配置。

MIPI BD D-PHY 由两个数据通道组成。配置 \hyperref[fielddesc:DSIHOSTNLANES]{DSI\_HOST\_N\_LANES} 以选择使用一个或两个数据通道。

默认情况下，PHY 设置为在 80-110 Mbps 的较低范围内操作。为了适应更高的比特率，可以使用适当的代码，在测试代码 \hyperref[sec:test-code-0x44]{测试代码 0x44} 中调整 HS 频率范围（hsfreqrange 字段）。如果 PHY 始终在相同的比特率下工作，建议在关闭状态下进行此调整，因为控制接口独立于 MIPI BD D-PHY 的其它部分。支持的范围如表 \ref{tab:mipidsi-hsfreq} 所示：

\begin{longtable}{| p{5cm} | R{5cm} |}
    \caption{MIPI BD D-PHY 支持的频率范围} \label{tab:mipidsi-hsfreq} \\
    \hline
    
    \rowcolor{lightgray}
    \textbf{频率范围 (Mbps)} & \textbf{hsfreqrange[5:0]\tnote{*}} \\ \hline
    \endfirsthead

    \multicolumn{2}{c}{\bfseries \tablename\ \thetable{} -- 接上页}\\\hline
    \rowcolor{lightgray}
    \textbf{频率范围 (Mbps)} & \textbf{hsfreqrange[5:0]\tnote{*}} \\ \hline
    \endhead

    \multicolumn{2}{r}{\textbf{见下页}} \\
    \endfoot

    \endlastfoot
    
    80-89 & 000000  \\ \hline
    90-99 & 010000  \\ \hline
    100-109 & 100000  \\ \hline
    110-129 & 000001  \\ \hline
    130-139 & 010001  \\ \hline
    140-149 & 100001  \\ \hline
    150-169 & 000010  \\ \hline
    170-179 & 010010  \\ \hline
    180-199 & 100010  \\ \hline
    200-219 & 000011  \\ \hline
    220-239 & 010011  \\ \hline
    240-249 & 100011  \\ \hline
    250-269 & 000100  \\ \hline
    270-299 & 010100  \\ \hline
    300-329 & 000101  \\ \hline
    330-359 & 010101  \\ \hline
    360-399 & 100101  \\ \hline
    400-449 & 000110  \\ \hline
    450-499 & 010110  \\ \hline
    500-549 & 000111  \\ \hline
    550-599 & 010111  \\ \hline
    600-649 & 001000  \\ \hline
    650-699 & 011000  \\ \hline
    700-749 & 001001  \\ \hline
    750-799 & 011001  \\ \hline
    800-849 & 101001  \\ \hline
    850-899 & 111001  \\ \hline
    900-949 & 001010  \\ \hline
    950-999 & 011010  \\ \hline
    1000-1049 & 101010  \\ \hline
    1050-1099 & 111010  \\ \hline
    1100-1149 & 001011  \\ \hline
    1150-1199 & 011011  \\ \hline
    1200-1249 & 101011  \\ \hline
    1250-1299 & 111011  \\ \hline
    1300-1349 & 001100  \\ \hline
    1350-1399 & 011100  \\ \hline
    1400-1449 & 101100  \\ \hline
    1450-1500 & 111100  \\ \hline
\end{longtable}

还需要通过测试代码接口配置 PLL 输入和输出频率。有关 PLL 参数的详细描述，请参阅 \hyperref[sec:pll-parameters]{PLL 参数}。

当 \hyperref[fielddesc:DSIHOSTPHYRSTZ]{DSI\_HOST\_PHY\_RSTZ} 和 \hyperref[fielddesc:DSIHOSTPHYSHUTDOWNZ]{DSI\_HOST\_PHY\_SHUTDOWNZ} 设置为 1 时，MIPI BD D-PHY 离开关闭状态并开始初始化。

\paragraph{PLL 锁定和 AFE 初始化}
在释放 \hyperref[fielddesc:DSIHOSTPHYRSTZ]{DSI\_HOST\_PHY\_RSTZ} 和 \hyperref[fielddesc:DSIHOSTPHYSHUTDOWNZ]{DSI\_HOST\_PHY\_SHUTDOWNZ} 后，PHY 开始初始化序列以启用其正常操作。建议在释放 \hyperref[fielddesc:DSIHOSTPHYRSTZ]{DSI\_HOST\_PHY\_RSTZ} 之前释放 \hyperref[fielddesc:DSIHOSTPHYSHUTDOWNZ]{DSI\_HOST\_PHY\_SHUTDOWNZ}。此外，确保 DSI\_DPHY\_CFG\_CLK 和 DSI\_DPHY\_PLL\_REFCLK（有关详细信息，请参见第 \ref{mod:resclk} 章中的表 \ref{tab:resclk-periph-clks-hp2} \textit{\nameref{mod:resclk}}）在此时可用且稳定。

之后，PLL 变为活动状态，PHY 等待 PLL 锁定。锁定状态由 \hyperref[fielddesc:DSIHOSTPHYLOCK]{DSI\_HOST\_PHY\_LOCK} 指示。

在 PLL 锁定后，开始 AFE 初始化。当 \hyperref[fielddesc:DSIHOSTPHYSTOPSTATECLKLANE]{DSI\_HOST\_PHY\_STOPSTATECLKLANE} \hyperref[fielddesc:DSIHOSTPHYSTOPSTATE0LANE]{DSI\_HOST\_PHY\_STOPSTATE\textcolor{red}{n}LANE} 被置位时，所有初始化步骤完成，通道处于停止状态 (LP-11)。

\paragraph{控制模式 (Control)}
AFE 初始化完成后，MIPI BD D-PHY 默认保持在控制模式，同时在线路上驱动停止状态 LP-11，直至有请求发出。任何请求都必须从该状态开始，并在该状态结束。请求发出后，通道可以进入高速数据传输模式、逃逸模式、超低功耗状态或转向状态。

\paragraph{高速数据传输}
DSI 主机在收到来自 DSI 桥接器的数据后，会发送高速数据传输请求。在高速数据传输模式下，物理层 (PHY) 以突发方式传输图像数据。通道仅在这些突发传输期间以高速模式运行。

如图 \ref{fig:mipidsi-hs} 所示，通道上的突发由三个部分组成：低功耗初始化序列、高速数据载荷和传输结束序列。

\begin{figure}[H]
     \centering
     \includegraphics[width=1.0\linewidth]{\modulefigures/mipidsi-hs.pdf}
     \caption{MIPI DSI HS 数据传输序列}
     \label{fig:mipidsi-hs}
 \end{figure}

\paragraph{逃逸模式 (Escape)}
逃逸模式允许在数据通道上以低速进行异步通信。在逃逸模式下，通道保持在低功耗状态。 \chipname{} 中的 MIPI DSI 在逃逸模式下支持低功耗数据传输 (LPDT) 和超低功耗状态 (ULPS)。

如果 \hyperref[fielddesc:DSIHOSTLPCMDEN]{DSI\_HOST\_LP\_CMD\_EN} 设置为 1，命令将在 LPDT 中发送。

当 DSI 空闲时，可以通过 \hyperref[regdesc:DSIHOSTPHYULPSCTRLREG]{DSI\_HOST\_PHY\_ULPS\_CTRL\_REG} 将 PHY 驱动到 ULPS 以节省功耗。通道状态可以通过 \hyperref[regdesc:DSIHOSTPHYSTATUSREG]{DSI\_HOST\_PHY\_STATUS\_REG} 观察得知。

\paragraph{转向}
PHY 可以通过转向交换传输方向。这用于接收来自 DSI 设备的命令响应。

\subsubsection{PLL 参数}
\label{sec:pll-parameters}
PLL 的 VCO 输出频率是输入参考时钟 DSI\_DPHY\_PLL\_REFCLK 的频率以及倍增比率 $M$ 和分频比率 $N$ 的函数：

\[
f_{VCO} = \frac{M}{N} \times f_{REFCLK}
\]

该函数必须满足以下限制：

\[
40\,\mathrm{MHz} \geq \frac{f_{REFCLK}}{N} \geq 5\,\mathrm{MHz}
\]

此外，$M$ 必须是偶数。$M$ 和 $N$ 的值可以通过 \hyperref[sec:test-code-0x17]{测试代码0x17} 和 \hyperref[sec:test-code-0x18]{测试代码0x18} 进行配置，配置值等于实际值减 1，如表 \ref{tab:m-n-ratios} 所示。之后，\hyperref[sec:test-code-0x19]{测试代码0x19} 的第 5 位和第 4 位应设置为 1，使 $M$ 和 $N$ 配置生效。

\begin{longtable}{|p{3cm}|p{2cm}|p{3cm}|p{2cm}|}
    \caption{M 和 N} \label{tab:m-n-ratios} \\
    \hline
    \rowcolor{lightgray}
    \textbf{m[8:0] = M-1} & \textbf{M} & \textbf{n[6:0] = N-1} & \textbf{N} \\ \hline
    \endfirsthead

    \multicolumn{4}{c}{\bfseries \tablename\ \thetable{} -- 接上页}\\\hline
    \rowcolor{lightgray}
    \textbf{m[8:0] = M-1} & \textbf{M} & \textbf{n[6:0] = N-1} & \textbf{N} \\ \hline
    \endhead

    \multicolumn{4}{r}{\textbf{见下页}} \\
    \endfoot

    \endlastfoot
    
    9'h1 & 2 & 7'h0 & 1  \\ \hline
    9'h3 & 4 & 7'h1 & 2  \\ \hline
    ... & ... & ... & ...  \\ \hline
\end{longtable}

VCO 的 $f_{VCO}$ 范围也可以通过 \hyperref[sec:test-code-0x10]{测试代码0x10} 进行选择。有效的 VCO 范围参见表 \ref{tab:vco-ranges}。

\begin{longtable}{|p{3cm}|p{5cm}|}
    \caption{VCO 的范围} \label{tab:vco-ranges} \\
    \hline
    \rowcolor{lightgray}
    \textbf{vcorange[2:0]} & \textbf{VCO 频率 (MHz)} \\ \hline
    \endfirsthead

    \multicolumn{2}{c}{\bfseries \tablename\ \thetable{} -- 接上页}\\\hline
    \rowcolor{lightgray}
    \textbf{vcorange[2:0]} & \textbf{VCO频率 (MHz)} \\ \hline
    \endhead

    \multicolumn{2}{r}{\textbf{见下页}} \\
    \endfoot

    \endlastfoot
    
    000 & 80-200  \\ \hline
    001 & 200-300  \\ \hline
    010 & 300-500  \\ \hline
    011 & 500-700  \\ \hline
    100 & 700-900  \\ \hline
    101 & 900-1100  \\ \hline
    110 & 1100-1300  \\ \hline
    111 & 1300-1500  \\ \hline
\end{longtable}

环路带宽应应根据输出频率通过测试代码接口进行配置，以确保 PLL 的正常运行。相关参数如下所示，有效值见表 \ref{tab:pll-cp-and-lpf-control-bits}。

\begin{itemize}
    \item VCO电容(vcocap[1:0])，通过 \hyperref[sec:test-code-0x10]{测试代码0x10} 配置。
    \item CP电流 (icpctrl[3:0])，通过 \hyperref[sec:test-code-0x11]{测试代码0x11} 配置。
    \item LPF特性 (lpfctrl[5:0])，通过 \hyperref[sec:test-code-0x12]{测试代码0x12} 配置。
\end{itemize}

\begin{longtable}{|p{3cm}|p{3cm}|p{3cm}|p{2cm}|p{2cm}|p{2cm}|}
    \caption{PLL CP 和 LPF 控制位} \label{tab:pll-cp-and-lpf-control-bits} \\
    \hline
    \rowcolor{lightgray}
    \textbf{VCO fmin (MHz)} & \textbf{VCO fmax (MHz)} & \textbf{vcorange[2:0]} & \textbf{icpctrl[3:0]} & \textbf{lpfctrl[5:0]} & \textbf{vcocap[1:0]} \\ \hline
    \endfirsthead

    \multicolumn{6}{c}{\bfseries \tablename\ \thetable{} -- 接上页}\\\hline
    \rowcolor{lightgray}
    \textbf{VCO fmin (MHz)} & \textbf{VCO fmax (MHz)} & \textbf{vcorange[2:0]} & \textbf{icpctrl[3:0]} & \textbf{lpfctrl[5:0]} & \textbf{vcocap[1:0]} \\ \hline
    \endhead

    \multicolumn{6}{r}{\textbf{见下页}} \\
    \endfoot

    \endlastfoot
    
    80   & 110  & 000 & 0100 & 010000 & 00 \\ \cline{1-2} \cline{4-6}
    110  & 150  &     & 1000 & 001000 & 00 \\ \cline{1-2} \cline{4-6}
    150  & 200  &     & 1100 & 001000 & 00 \\ \hline
    200  & 250  & 001 & 1100 & 001000 & 00 \\ \cline{1-2} \cline{4-6}
    250  & 300  &     & 0010 & 001000 & 00 \\ \hline
    300  & 400  & 010 & 0101 & 001000 & 00 \\ \cline{1-2} \cline{4-6}
    400  & 500  &     & 0101 & 000001 & 00 \\ \hline
    500  & 600  & 011 & 0110 & 010000 & 00 \\ \cline{1-2} \cline{4-6}
    600  & 700  &     & 0110 & 000100 & 00 \\ \hline
    700  & 900  & 100 & 0110 & 000100 & 00 \\ \hline
    900  & 1100 & 101 & 0111 & 010000 & 00 \\ \hline
    1100 & 1300 & 110 & 0111 & 001000 & 00 \\ \hline
    1300 & 1500 & 111 & 0111 & 001000 & 00 \\ \hline
\end{longtable}

\begin{tiplisting}
    vcorange, icpctrl, lpfctrl 在对应的 hsfreqrange 为 0 时会自动设置为默认值。
    \begin{itemize}
        \item \hyperref[sec:test-code-0x10]{测试代码 0x10，第 7 位}，清 0 以绕过 vcorange
        \item \hyperref[sec:test-code-0x12]{测试代码 0x12，第 7 位}，清 0 以绕过 icpctrl
        \item \hyperref[sec:test-code-0x12]{测试代码 0x12，第 6 位}，清 0 以绕过 lpfctrl
    \end{itemize}
\end{tiplisting}

\subsubsection{测试代码}
MIPI BD D-PHY 提供了一组测试代码，用于配置 MIPI BD D-PHY 操作。配置可以使用 \hyperref[regdesc:DSIHOSTPHYTSTCTRL0REG]{DSI\_HOST\_PHY\_TST\_CTRL0\_REG} 和 \hyperref[regdesc:DSIHOSTPHYTSTCTRL1REG]{DSI\_HOST\_PHY\_TST\_CTRL1\_REG} 通过测试代码接口完成。有关测试代码接口的详细信息，请参阅 \ref{mod:mipicsi} \textit{\nameref{mod:mipicsi}} 章 \ref{sec:mipicsi-rx-dphy-program-interface} 小节。

支持的所有测试代码包括：

\paragraph{正常运行}
\label{sec:normal-operation}
此模式表示 PHY 的正常运行以及测试接口的唤醒状态。在此模式下，测试接口处于非活动状态。

\begin{itemize}
    \item \textbf{测试代码}：0x00
    \item \textbf{测试数据}：无
\end{itemize}

\paragraph{PLL 偏置电流选择器 / 滤波电容控制 / VCO 控制}
\label{sec:test-code-0x10}
控制 PLL 的偏置和 VCO 的传递函数。

\begin{itemize}
    \item \textbf{测试代码}：0x10
    \item \textbf{测试数据}
        \begin{longtable}{ |p{3cm}|p{3cm}|p{3cm}|p{3cm}|p{3cm}|}
            \caption{0x10 的测试数据}
            \label{tab:test-data-0x10} \\
            \hline
            \rowcolor{lightgray}
            \textbf{w-1'b0} & \textbf{w-1'b0} & \textbf{w-3'b000} & \textbf{w-2'b00} & \textbf{w-1'b1} \\ \hline
            跳过 VCO 范围 & 保留 & VCO 范围控制 & VCO 内部电容控制 & 参考偏置电流选择器  \\ \hline
        \end{longtable}
        \begin{itemize}
            \item 位 7: 跳过 VCO 范围
                \begin{itemize}
                    \item 0 - VCO 范围使用对应 hsfreqrange 的默认值进行配置
                    \item 1 - VCO 范围使用位 5-3 进行配置
                \end{itemize}
            \item 位 6: 保留
            \item 位 5-3: VCO 范围控制，详细信息请参见表 \ref{tab:vco-ranges}
            \item 位 2-1: VCO 内部电容控制
                \begin{itemize}
                    \item 00 - 默认电容
                    \item 01 - 低电容（1/4 默认值）
                    \item 10 - 高电容（2 倍默认值）
                    \item 11 - 不允许
                \end{itemize}
            \item 位 0: 参考偏置电流选择器
                \begin{itemize}
                    \item 0 - 内部偏置
                    \item 1 -  20uA 清洁的外部电流源
                \end{itemize}
        \end{itemize}
    \item \textbf{测试数据输出}
        \begin{longtable}{ |p{3cm}|p{3cm}|p{3cm}|p{3cm}|p{3cm}|}
            \caption{0x10 的测试数据输出}
            \label{tab:test-dout-0x10} \\
            \hline
            \rowcolor{lightgray}
            \textbf{r-1'b0} & \textbf{r-1'b0} & \textbf{r-3'b000} & \textbf{r-2'b00} & \textbf{r-1'b1} \\ \hline
            跳过 VCO 范围回环 & 1'b0 & VCO 范围控制回环 & VCO 内部电容控制回环 & 参考偏置电流选择器回环  \\ \hline
        \end{longtable}
\end{itemize}

\paragraph{用于初始化和超低功耗模式的 PLL CP 控制/PLL 锁定旁路}
\label{sec:test-code-0x11}
控制充电泵 (CP) 电流和 PLL 锁定旁路，用于初始化和 ULP 退出。

\begin{itemize}
    \item \textbf{测试代码}：0x11
    \item \textbf{测试数据}
        \begin{longtable}{ |p{6cm}|p{6cm}|}
            \caption{0x11 的测试数据}
            \label{tab:test-data-0x11} \\
            \hline
            \rowcolor{lightgray}
            \textbf{w-4'b0000} & \textbf{w-4'b0010} \\ \hline
            PLL 锁定跳过 & CP 电流  \\ \hline
        \end{longtable}
        \begin{itemize}
            \item 位 7-4: PLL 锁定跳过
                \begin{itemize}
                    \item xxx1 - 在初始化期间跳过 PLL 锁定
                    \item x11x - 在退出 ULP 后跳过 PLL 锁定
                \end{itemize}
            \item 位 3-0: CP 电流 (uA)
                \begin{itemize}
                    \item 0000 - 1.5
                    \item 0001 - 3
                    \item 0010 - 4.5
                    \item 0011 - 6
                    \item 0100 - 2.5
                    \item 0101 - 5
                    \item 0110 - 7.5
                    \item 0111 - 10
                    \item 1000 - 3
                    \item 1001 - 6
                    \item 1010 - 9
                    \item 1011 - 12
                    \item 1100 - 4
                    \item 1101 - 8
                    \item 1110 - 12
                    \item 1111 - 16
                \end{itemize}
        \end{itemize}
    \item \textbf{测试数据输出}
        \begin{longtable}{ |p{6cm}|p{6cm}|}
            \caption{0x11 的测试数据输出}
            \label{tab:test-dout-0x11} \\
            \hline
            \rowcolor{lightgray}
            \textbf{r-1'b0} & \textbf{r-1'b0} \\ \hline
            PLL 锁定跳过回环 & CP 电流回环 \\ \hline
        \end{longtable}
\end{itemize}

\paragraph{PLL LPF 和 CP 控制}
\label{sec:test-code-0x12}
控制充电泵 (CP) 电流和低通滤波器 (LPF) 电阻。

\begin{itemize}
    \item \textbf{测试代码}：0x12
    \item \textbf{测试数据}
        \begin{longtable}{ |p{4cm}|p{4cm}|p{4cm}|}
            \caption{0x12 的测试数据}
            \label{tab:test-data-0x12} \\
            \hline
            \rowcolor{lightgray}
            \textbf{w-1'b0} & \textbf{w-1'b0} & \textbf{w-6'b000010} \\ \hline
            跳过 CP 默认值 & 跳过 LPF 默认值 & LPF 电阻  \\ \hline
        \end{longtable}
        \begin{itemize}
            \item 位 7: 跳过 CP 默认值
                \begin{itemize}
                    \item 0 - CP 使用对应 hsfreqrange 的默认值进行配置
                    \item 1 - CP 使用测试代码 0x11 的位 3-0 进行配置
                \end{itemize}
            \item 位 6: 跳过 LPF 默认值
                \begin{itemize}
                    \item 0 - LPF 使用对应 hsfreqrange 的默认值进行配置
                    \item 1 - LPF 使用位 5-0 进行配置
                \end{itemize}
            \item 位 5-0: LPF 电阻 ($k\Omega$)
                \begin{itemize}
                    \item 000000 - 18
                    \item 000001 - 15.6
                    \item 000010 - 15
                    \item 000100 - 14.4
                    \item 001000 - 12.8
                    \item 010000 - 11.4
                    \item 100000 - 10.5
                    \item 其他 - 保留
                \end{itemize}
        \end{itemize}
    \item \textbf{测试数据输出}
        \begin{longtable}{ |p{4cm}|p{4cm}|p{4cm}|}
            \caption{0x12 的测试数据输出}
            \label{tab:test-dout-0x12} \\
            \hline
            \rowcolor{lightgray}
            \textbf{r-1'b0} & \textbf{r-1'b0} & \textbf{r-6'b000010} \\ \hline
            CP 跳过回环 & LPF 跳过回环 & LPF 电阻回环 \\ \hline
        \end{longtable}
\end{itemize}


\paragraph{PLL 输入分频比}
\label{sec:test-code-0x17}
控制 PLL 的输入分频器 $N$。$N = n + 1$，其中 $N$ 是实际的分频比，$n$ 是测试代码中的配置值。

要启用这一配置值，请确保 \hyperref[sec:test-code-0x19]{测试代码 0x19} 的位 4 设置为 1。

\begin{itemize}
    \item \textbf{测试代码}：0x17
    \item \textbf{测试数据}
        \begin{longtable}{ |p{6cm}|p{6cm}|}
            \caption{0x17 的测试数据}
            \label{tab:test-data-0x17} \\
            \hline
            \rowcolor{lightgray}
            \textbf{w-1'b0} & \textbf{w-7'b0000001} \\ \hline
            保留 & 输入分频器  \\ \hline
        \end{longtable}
    \item \textbf{测试数据输出}
        \begin{longtable}{ |p{6cm}|p{6cm}|}
            \caption{0x17 的测试数据输出}
            \label{tab:test-dout-0x17} \\
            \hline
            \rowcolor{lightgray}
            \textbf{r-1'b0} & \textbf{r-7'b0000001} \\ \hline
            保留 & 输入分频器回环 \\ \hline
        \end{longtable}
\end{itemize}

\paragraph{PLL 循环分频比}
\label{sec:test-code-0x18}
控制 PLL 的反馈分频器 $M$。$M = m + 1$，其中 $M$ 是实际的反馈比，$m$ 是测试代码中的配置值。

要启用配置值，请确保 \hyperref[sec:test-code-0x19]{测试代码 0x19} 的位 5 设置为 1。

\begin{itemize}
    \item \textbf{测试代码}：0x18
    \item \textbf{测试数据}
        \begin{longtable}{ |p{4cm}|p{4cm}|p{4cm}|}
            \caption{0x18 的测试数据}
            \label{tab:test-data-0x18} \\
            \hline
            \rowcolor{lightgray}
            \textbf{w-1'b0} & \textbf{w-2'b00} & \textbf{w-5'b00101} \\ \hline
            设置 0 & 保留 & 反馈分频器 [4:0]  \\ \hline
            \rowcolor{lightgray}
            \textbf{w-1'b1} & \textbf{w-3'b00} & \textbf{w-4'b0000} \\ \hline
            设置 1 & 保留 & 反馈分频器 [8:5]  \\ \hline
        \end{longtable}
        \begin{itemize}
            \item 位 7: 配置选择器
                \begin{itemize}
                    \item 0 - 设置 0 被配置
                    \item 1 - 设置 1 被配置
                \end{itemize}
            \item 位 6-5: 保留
            \item 位 4-0: 反馈分频器
        \end{itemize}
    \item \textbf{测试数据输出}\\
        TESTDIN[7] - 设置 0
        \begin{longtable}{ |p{4cm}|p{8cm}|}
            \caption{0x18 的测试数据输出}
            \label{tab:test-dout-0x18-set0} \\
            \hline
            \rowcolor{lightgray}
            \textbf{r-3'b000} & \textbf{r-5'b00101} \\ \hline
            3'b000 & 反馈分频器 [4:0] 回环 \\ \hline
        \end{longtable}
        TESTDIN[7] - 设置 1
        \begin{longtable}{ |p{4cm}|p{8cm}|}
            \caption{0x18 的测试数据输出}
            \label{tab:test-dout-0x18-set1} \\
            \hline
            \rowcolor{lightgray}
            \textbf{r-4'b0000} & \textbf{r-4'b0000} \\ \hline
            4'b0000 & 反馈分频器 [8:5] 回环 \\ \hline
        \end{longtable}
\end{itemize}

\paragraph{PLL 输入和循环分频比控制}
\label{sec:test-code-0x19}
控制 PLL 输入和循环分频比。

\begin{itemize}
    \item \textbf{测试代码}：0x19
    \item \textbf{测试数据}
        \begin{longtable}{ |p{2cm}|p{5cm}|p{5cm}|p{2cm}|}
            \caption{0x19 的测试数据}
            \label{tab:test-data-0x19} \\
            \hline
            \rowcolor{lightgray}
            \textbf{w-2'b00} & \textbf{w-1'b0} & \textbf{w-1'b0} & \textbf{w-4'b0000} \\ \hline
            保留 & 跳过 PLL 循环分频器默认值 & 跳过 PLL 输入分频器默认值 & 保留 \\ \hline
        \end{longtable}
        \begin{itemize}
            \item 位 7-6: 保留
            \item 位 5: 跳过 PLL 循环分频器默认值
                \begin{itemize}
                    \item 0 - PLL 循环分频器使用对应 hsfreqrange 的默认值进行配置
                    \item 1 - PLL 循环分频器使用 \hyperref[sec:test-code-0x18]{测试代码 0x18} 进行配置
                \end{itemize}
            \item 位 4: 跳过 PLL 输入分频器默认值
                \begin{itemize}
                    \item 0 - PLL 输入分频器使用对应 hsfreqrange 的默认值进行配置
                    \item 1 - PLL 输入分频器使用 \hyperref[sec:test-code-0x17]{测试代码 0x17} 进行配置
                \end{itemize}
            \item 位 3-0: 保留
        \end{itemize}
    \item \textbf{测试数据输出}
        \begin{longtable}{ |p{2cm}|p{5cm}|p{5cm}|p{2cm}|}
            \caption{0x19 的测试数据输出}
            \label{tab:test-dout-0x19} \\
            \hline
            \rowcolor{lightgray}
            \textbf{r-2'b00} & \textbf{r-1'b0} & \textbf{r-1'b0} & \textbf{r-4'b0000} \\ \hline
            保留 & 跳过 PLL 循环分频器默认回环 & 跳过 PLL 输入分频器默认回环 & 保留 \\ \hline
        \end{longtable}
\end{itemize}

\paragraph{AFE/BIAS/Bandgap 模拟可编程性}
\label{sec:test-code-0x22}
控制 PHY 中的各种参数。

\begin{itemize}
    \item \textbf{测试代码}：0x22
    \item \textbf{测试数据}：
        \begin{longtable}{ |p{2cm}|p{4cm}|p{6cm}| }
            \caption{0x22 的测试数据}
            \label{tab:test-data-0x22} \\
            \hline
            \rowcolor{lightgray}
            \textbf{w-1'b0} & \multicolumn{2}{l|}{\textbf{w-7'b0000000}} \\ \hline
            设置 0 & \multicolumn{2}{l|}{模拟可编程性 [6:0]}  \\ \hline
            \rowcolor{lightgray}
            \textbf{w-1'b1} & \textbf{w-3'b00} & \textbf{w-4'b0000} \\ \hline
            设置 1 & 保留 & 模拟可编程性 [10:7]  \\ \hline
        \end{longtable}

        \begin{itemize}
            \item 位 7：配置选择器
                \begin{itemize}
                    \item 0 - 配置设置 0（同时选择位 6:0 用于可观测性）
                    \item 1 - 配置设置 1（同时选择位 10:7 用于可观测性）
                \end{itemize}
            \item 位 6-0：模拟可编程性
                \begin{itemize}
                    \item 位 10：1.5 Gbps 支持。对于高于 1 Gbps 且不超过 1.5 Gbps 的速率，应启用此位
                        \begin{itemize}
                            \item 0 - 禁用 1.5 Gbps 模拟电路支持
                            \item 1 - 启用 1.5 Gbps 模拟电路支持
                        \end{itemize}
                    \item 位 9-7：调整带隙参考电压（典型角）
                        \begin{itemize}
                            \item 000 - 97.07\%
                            \item 001 - 98.05\%
                            \item 010 - 99.02\%
                            \item 011 - 100.00\%
                            \item 100 - 93.17\%
                            \item 101 - 94.15\%
                            \item 110 - 95.12\%
                            \item 111 - 96.10\%
                        \end{itemize}
                    \item 位 6：调整 LP RX 偏置电流
                        \begin{itemize}
                            \item 0 - 100\%
                            \item 1 - 50\%
                        \end{itemize}
                    \item 位 5-3：调整 LP TX 偏置电流
                        \begin{itemize}
                            \item 000 - 76.92\%
                            \item 001 - 83.33\%
                            \item 010 - 90.91\%
                            \item 011 - 100.00\%
                            \item 100 - 111.10\%
                            \item 101 - 125.00\%
                            \item 110 - 142.80\%
                            \item 111 - 166.70\%
                        \end{itemize}
                    \item 位 2-0：Biasextr：调整内部电阻
                        \begin{itemize}
                            \item 000 - 87.1\%
                            \item 001 - 91.5\%
                            \item 010 - 95.9\%
                            \item 011 - 100\%
                            \item 100 - 105.94\%
                            \item 101 - 111.88\%
                            \item 110 - 118.8\%
                            \item 111 - 127.7\%
                        \end{itemize}
                \end{itemize}
        \end{itemize}
    \item \textbf{测试数据输出}：\\
        TESTDIN[7] - 设置 0
        \begin{longtable}{ |p{3cm}|p{9cm}|}
            \caption{0x22 的测试数据输出}
            \label{tab:test-dout-0x22-set0} \\
            \hline
            \rowcolor{lightgray}
            \textbf{r-1'b0} & \textbf{r-7'b0000000} \\ \hline
            1'b0 & 模拟可编程性 [6:0] 回环 \\ \hline
        \end{longtable}
        TESTDIN[7] - 设置 1
        \begin{longtable}{ |p{3cm}|p{9cm}|}
            \caption{0x22 的测试数据输出}
            \label{tab:test-dout-0x22-set1} \\
            \hline
            \rowcolor{lightgray}
            \textbf{r-4'b0000} & \textbf{r-4'b0000} \\ \hline
            4'b0000 & 模拟可编程性 [10:7] 回环 \\ \hline
        \end{longtable}
\end{itemize}

\paragraph{HS RX 控制通道 0}
\label{sec:test-code-0x44}
控制 hsfreqrange。

\begin{itemize}
    \item \textbf{测试代码}：0x44
    \item \textbf{测试数据}
        \begin{longtable}{ |p{2cm}|p{6cm}|p{6cm}|}
            \caption{0x44 的测试数据}
            \label{tab:test-data-0x44} \\
            \hline
            \rowcolor{lightgray}
            \textbf{w-1'b0} & \textbf{w-6'b000000} & \textbf{w-1'b0} \\ \hline
            设置 2 & HS 操作频率范围选择 (hsfreqrange) & 跳过 HS-RX 稳定滤波器 \\ \hline
        \end{longtable}

        \begin{itemize}
            \item 位 7: 配置选择器
                \begin{itemize}
                    \item 0 - 设置 2 被配置
                    \item 1 - 无效
                \end{itemize}
            \item 位 6-1: HS 频率范围选择。详细值请参见表 \ref{tab:mipidsi-hsfreq}。
            \item 位 0: 跳过 HS 稳定滤波器
        \end{itemize}
    \item \textbf{测试数据输出}
        \begin{longtable}{ |p{2cm}|p{6cm}|p{6cm}|}
            \caption{0x44 的测试数据输出}
            \label{tab:test-dout-0x44} \\
            \hline
            \rowcolor{lightgray}
            \textbf{r-1'b0} & \textbf{r-6'b000000} & \textbf{r-1'b0} \\ \hline
            1'b0 & HS 运行的频率范围选择回环 (hsfreqrange) & 跳过 HS-RX 稳定滤波器回环 \\ \hline
        \end{longtable}
\end{itemize}

% 
\subsection{时钟和显示分辨率}
\label{sec:clocks-and-display-resolution}

DSI 的时钟架构如图 \ref{figure:mipidsi-clk-arch} 所示。所有时钟为：
\begin{itemize}
    \item \textbf{PCLK}：APB 总线时钟
    \item \textbf{ACLK}：AXI 总线时钟
    \item \textbf{DSI\_DPICLK}：DPI 接口时钟
    \item \textbf{DSI\_DPHY\_CFG\_CLK}：DPHY 配置时钟
    \item \textbf{DSI\_DPHY\_PLL\_REFCLK}：DPHY PLL 输入参考时钟
    \item \textbf{TX\_HS\_CLK}：DPHY PLL 输出时钟
    \item \textbf{TX\_BYTE\_CLK}：高速 (HS) 传输字节时钟，TX\_HS\_CLK 除以 8。
    \item \textbf{TX\_ESC\_CLK}：低功耗 (LP) 逃逸模式传输时钟，从 TX\_BYTE\_CLK 分出。分频比可以通过 \hyperref[fielddesc:DSIHOSTTOCLKDIVISION]{DSI\_HOST\_TO\_CLK\_DIVISION} 配置，允许的最大频率为 20 MHz。
\end{itemize}

\begin{figure}[H]
    \centering
    \includegraphics[width=1.0\textwidth]{\modulefigures/mipidsi-clk-arch.pdf}
    \caption{MIPI DSI 时钟架构}
    \label{figure:mipidsi-clk-arch}
\end{figure}

DSI 支持连续和非连续时钟通道。
\begin{itemize}
    \item 对于连续时钟通道，需要把 \hyperref[fielddesc:DSIHOSTAUTOCLKLANECTRL]{DSI\_HOST\_AUTO\_CLKLANE\_CTRL} 配置成 1, 把 \hyperref[fielddesc:DSIHOSTPHYTXREQUESTCLKHS]{DSI\_HOST\_PHY\_TXREQUESTCLKHS} 配置成 1
    \item 对于非连续时钟通道，需要把 \hyperref[fielddesc:DSIHOSTAUTOCLKLANECTRL]{DSI\_HOST\_AUTO\_CLKLANE\_CTRL} 配置成 0, 把 \hyperref[fielddesc:DSIHOSTPHYTXREQUESTCLKHS]{DSI\_HOST\_PHY\_TXREQUESTCLKHS} 配置成 1
\end{itemize}

在视频模式传输中，视频时序参数如图 \ref{figure:mipidsi-video-timing} 所示。
\begin{figure}[H]
    \centering
    \includegraphics[width=0.8\textwidth]{\modulefigures/mipidsi-video-timing.pdf}
    \caption{视频时序}
    \label{figure:mipidsi-video-timing}
\end{figure}

对于 DPI，视频时序和分辨率可以通过以下 DSI 桥接寄存器字段进行配置。之后，向 \hyperref[fielddesc:DSIBRGDPICONFIGUPDATE]{DSI\_BRG\_DPI\_CONFIG\_UPDATE} 写入 1 使配置生效。
\begin{itemize}
    \item \hyperref[fielddesc:DSIBRGHTOTAL]{DSI\_BRG\_HTOTAL}：以 DSI\_DPICLK 为单位的总水平长度 (HTOTAL)
    \item \hyperref[fielddesc:DSIBRGHDISP]{DSI\_BRG\_HDISP}：以 DSI\_DPICLK 为单位的有效水平长度 (HACT)
    \item \hyperref[fielddesc:DSIBRGHBANK]{DSI\_BRG\_HBANK}：以 DSI\_DPICLK 为单位的水平后沿长度 (HBP)
    \item \hyperref[fielddesc:DSIBRGHSYNC]{DSI\_BRG\_HSYNC}：以 DSI\_DPICLK 为单位的水平同步长度 (HSA)
    \item \hyperref[fielddesc:DSIBRGVTOTAL]{DSI\_BRG\_VTOTAL}：以行数为单位的总垂直长度 (VTOTAL)
    \item \hyperref[fielddesc:DSIBRGVDISP]{DSI\_BRG\_VDISP}：以行数为单位的有效垂直长度(VACT)
    \item \hyperref[fielddesc:DSIBRGVBANK]{DSI\_BRG\_VBANK}：以行数为单位的垂直后沿长度 (VBP)
    \item \hyperref[fielddesc:DSIBRGVSYNC]{DSI\_BRG\_VSYNC}：以行数为单位的垂直同步长度 (VSA)

\end{itemize}


DSI 也需要类似的视频时序和分辨率配置，相关寄存器字段如下：
\begin{itemize}
    \item \hyperref[fielddesc:DSIHOSTVIDHSATIME]{DSI\_HOST\_VID\_HSA\_TIME}：以 TX\_BYTE\_CLK 为单位的水平同步长度 (HSA)
    \item \hyperref[fielddesc:DSIHOSTVIDHBPTIME]{DSI\_HOST\_VID\_HBP\_TIME}：以 TX\_BYTE\_CLK 为单位的水平后沿长度 (HBP)
    \item \hyperref[fielddesc:DSIHOSTVIDHLINETIME]{DSI\_HOST\_VID\_HLINE\_TIME}：以 TX\_BYTE\_CLK 为单位的总水平长度 (HTOTAL)
    \item \hyperref[fielddesc:DSIHOSTVSALINES]{DSI\_HOST\_VSA\_LINES}：以行数为单位的垂直同步长度 (VSA)
    \item \hyperref[fielddesc:DSIHOSTVBPLINES]{DSI\_HOST\_VBP\_LINES}：以行数为单位的垂直后沿长度 (VBP)
    \item \hyperref[fielddesc:DSIHOSTVFPLINES]{DSI\_HOST\_VFP\_LINES}：以行数为单位的垂直前沿长度 (VFP)
    \item \hyperref[fielddesc:DSIHOSTVACTIVELINES]{DSI\_HOST\_VACTIVE\_LINES}：以行数为单位的有效垂直长度 (VACT)
\end{itemize}

DPI 和 DSI 的视频时序配置应匹配，以避免行缓冲区溢出。DPI 和 DSI 中每行的实际时序应相同。

%%%%%%%%%%%%%%%%%%%%%
%%% 中断 %%%
%%%%%%%%%%%%%%%%%%%%%

\section{中断}
\chipname{} 的 MIPI DSI 可以生成以下中断信号，这些信号将发送到 \hyperref[mod:intmtrx]{中断矩阵}。

\begin{itemize}
    \item dsi\_bridge\_interrupt
    \item dsi\_interrupt
\end{itemize}

以下中断源可以生成 dsi\_bridge\_interrupt 中断信号：

\begin{longtable}[c]{ |p{5cm}|p{10cm}| }
    \caption{DSI 桥中断}
    \label{tab:dsi-bridge-intr}\\\hline

    \rowcolor{lightgray}
    \textbf{中断源} & \textbf{详细信息} \\\hline
    \endfirsthead


    \multicolumn{2}{c}{\bfseries \tablename\ \thetable{} -- ...接上页}\\\hline
    \rowcolor{lightgray}
    \textbf{中断源} & \textbf{详细信息} \\\hline
    \endhead

    \multicolumn{2}{r}{\textbf{见下页...}} \\
    \endfoot

    %\insertTableNotes
    \endlastfoot

    \hyperref[fielddesc:DSIBRGUNDERRUNINTRAW]{DSI\_BRG\_UNDERRUN\_INT} &  DSI 桥缓冲下溢时触发 \\\hline
        

\end{longtable}

以下中断源可以生成 dsi\_interrupt 中断信号：

\begin{longtable}[c]{ |p{5cm}|p{5cm}|p{5cm}| }
    \caption{DSI 主机中断}
    \label{tab:dsi-host-interrupt}\\\hline

    \rowcolor{lightgray}
    \textbf{名称} & \textbf{原因} & \textbf{推荐处理方法}  \\\hline
    \endfirsthead


    \multicolumn{3}{c}{\bfseries \tablename\ \thetable{} -- ...接上页}\\\hline
    \rowcolor{lightgray}
    \textbf{名称} & \textbf{原因} & \textbf{推荐处理方法}  \\\hline
    \endhead

    \multicolumn{3}{r}{\textbf{见下页...}} \\
    \endfoot

    %\insertTableNotes
    \endlastfoot

    \hyperref[fielddesc:DSIHOSTACKWITHERR0]{\makecell[l]{DSI\_HOST\_ACK\_WITH\_ERR\_0}} & 确认错误报告中的 SoT 错误 & 如果数据包很关键，请重新传输数据包 \\\hline
    \hyperref[fielddesc:DSIHOSTACKWITHERR1]{\makecell[l]{DSI\_HOST\_ACK\_WITH\_ERR\_1}} & 确认错误报告中的 SoT 同步错误 & 请重新传输数据包 \\\hline
    \hyperref[fielddesc:DSIHOSTACKWITHERR2]{\makecell[l]{DSI\_HOST\_ACK\_WITH\_ERR\_2}} & 确认错误报告中的 EoT 同步错误 & 请重新传输数据包 \\\hline
    \hyperref[fielddesc:DSIHOSTACKWITHERR3]{\makecell[l]{DSI\_HOST\_ACK\_WITH\_ERR\_3}} & 确认错误报告中的转义模式入口命令错误，设备未识别转义模式入口代码 & 检查设备能力 \\\hline
    \hyperref[fielddesc:DSIHOSTACKWITHERR4]{\makecell[l]{DSI\_HOST\_ACK\_WITH\_ERR\_4}} & 确认错误报告中的低功耗传输同步错误 & 请重新传输数据包 \\\hline
    \hyperref[fielddesc:DSIHOSTACKWITHERR5]{\makecell[l]{DSI\_HOST\_ACK\_WITH\_ERR\_5}} & 确认错误报告中的外设超时错误 & 检查主机超时计数器配置是否与设备规格一致；主机超时计数器配置应比设备更长。此错误会自动恢复，但不能保证传输或接收完成\\\hline
    \hyperref[fielddesc:DSIHOSTACKWITHERR6]{\makecell[l]{DSI\_HOST\_ACK\_WITH\_ERR\_6}} & 确认错误报告中的虚假控制错误 & 重置 DSI 主机和 DPHY \\\hline
    \hyperref[fielddesc:DSIHOSTACKWITHERR7]{\makecell[l]{DSI\_HOST\_ACK\_WITH\_ERR\_7}} & 确认错误报告中的保留项（特定于设备） & 检查设备能力和连接性并重新传输 \\\hline
    \hyperref[fielddesc:DSIHOSTACKWITHERR8]{\makecell[l]{DSI\_HOST\_ACK\_WITH\_ERR\_8}} & 确认错误报告中的 ECC 错误，单比特（已检测并纠正） & 无需采取行动。如果此错误频繁发生，请分析信号完整性或链路的噪声条件 \\\hline
    \hyperref[fielddesc:DSIHOSTACKWITHERR9]{\makecell[l]{DSI\_HOST\_ACK\_WITH\_ERR\_9}} & 确认错误报告中的 ECC 错误，多比特（已检测未纠正） & 如果数据包很关键，请重新传输数据包。如果此错误频繁发生，请分析信号完整性或链路的噪声条件\\\hline
    \hyperref[fielddesc:DSIHOSTACKWITHERR10]{\makecell[l]{DSI\_HOST\_ACK\_WITH\_ERR\_10}} & 确认错误报告中的校验和错误（仅限长数据包） & 如果数据包很关键，请重新传输数据包。如果此错误频繁发生，请检查连接 \\\hline
    \hyperref[fielddesc:DSIHOSTACKWITHERR11]{\makecell[l]{DSI\_HOST\_ACK\_WITH\_ERR\_11}} & 确认错误报告中的未识别 DSI 数据类型。某些数据包可能不受支持 & 检查设备能力 \\\hline
    \hyperref[fielddesc:DSIHOSTACKWITHERR13]{\makecell[l]{DSI\_HOST\_ACK\_WITH\_ERR\_13}} & 确认错误报告中的无效传输长度 & 请重新传输数据包。如果此错误频繁发生，请检查连接 \\\hline
    \hyperref[fielddesc:DSIHOSTACKWITHERR14]{\makecell[l]{DSI\_HOST\_ACK\_WITH\_ERR\_14}} & 确认错误报告中的保留项（特定于设备） & 检查设备规格以找出可能的原因和措施\\\hline
    \hyperref[fielddesc:DSIHOSTACKWITHERR15]{\makecell[l]{DSI\_HOST\_ACK\_WITH\_ERR\_15}} & 确认错误报告中的 DSI 协议违规 & 检查显示规格。发生此错误时，设备应有另一个读回命令，报告附加错误信息 \\\hline
    \hyperref[fielddesc:DSIHOSTDPHYERRORS0]{\makecell[l]{DSI\_HOST\_DPHY\_ERRORS\_0}} & 来自通道 0 的 ErrEsc 转义入口错误，未识别转义入口代码 & 重复接收 \\\hline
    \hyperref[fielddesc:DSIHOSTDPHYERRORS1]{\makecell[l]{DSI\_HOST\_DPHY\_ERRORS\_1}} & 来自通道 0 的 ErrSyncEsc 低功耗数据传输同步错误 & 重置 DSI 主机并重复传输 \\\hline
    \hyperref[fielddesc:DSIHOSTDPHYERRORS2]{\makecell[l]{DSI\_HOST\_DPHY\_ERRORS\_2}} & 来自通道 0 的控制错误 ErrControl & 重置 DSI 主机和 DPHY \\\hline
    \hyperref[fielddesc:DSIHOSTDPHYERRORS3]{\makecell[l]{DSI\_HOST\_DPHY\_ERRORS\_3}} & 来自通道 0 的 LP0 争用错误 ErrContentionLP0 & 重置 DSI 主机并重复传输 \\\hline
    \hyperref[fielddesc:DSIHOSTDPHYERRORS4]{\makecell[l]{DSI\_HOST\_DPHY\_ERRORS\_4}} & 来自通道 0 的 LP1 争用错误 ErrContentionLP1 & 重置 DSI 主机并重复传输 \\\hline
    \hyperref[fielddesc:DSIHOSTTOHSTX]{\makecell[l]{DSI\_HOST\_TO\_HS\_TX}} & 高速传输超时错误 & 检查超时配置。主机超时计数器配置应比设备更长 \\\hline
    \hyperref[fielddesc:DSIHOSTTOLPRX]{\makecell[l]{DSI\_HOST\_TO\_LP\_RX}} & 低功耗接收超时错误 & 检查超时配置。主机超时计数器配置应比设备更长 \\\hline
    \hyperref[fielddesc:DSIHOSTECCSINGLEERR]{\makecell[l]{DSI\_HOST\_ECC\_SINGLE\_ERR}} & 在接收的数据包中检测到并纠正了 ECC 单个错误 & 无需采取任何行动 \\\hline
    \hyperref[fielddesc:DSIHOSTECCMILTIERR]{\makecell[l]{DSI\_HOST\_ECC\_MILTI\_ERR}} & 在接收的数据包中检测到 ECC 多个错误 & 重复传输 \\\hline
    \hyperref[fielddesc:DSIHOSTCRCERR]{\makecell[l]{DSI\_HOST\_CRC\_ERR}} & 在接收的数据包有效载荷中检测到 CRC 错误 & 重置 DSI 主机并重复传输 \\\hline
    \hyperref[fielddesc:DSIHOSTPKTSIZEERR]{\makecell[l]{DSI\_HOST\_PKT\_SIZE\_ERR}} & 在数据包接收期间检测到数据包大小错误 & 重置 DSI 主机并重复传输 \\\hline
    \hyperref[fielddesc:DSIHOSTEOPTERR]{\makecell[l]{DSI\_HOST\_EOPT\_ERR}} & 在外设传输结束时未接收到 EoTp 数据包 & 此错误对接收数据包的数据完整性没有关键影响。检查设备是否支持 EoTp 数据包 \\\hline
    \hyperref[fielddesc:DSIHOSTDPIPLDWRERR]{\makecell[l]{DSI\_HOST\_DPI\_PLD\_WR\_ERR}} & 在 DPI 像素行存储期间，有效载荷 FIFO 变满且存储的数据已损坏 & 检查视频模式配置，重置 DSI 主机并重新启动视频传输 \\\hline
    \hyperref[fielddesc:DSIHOSTGENCMDWRERR]{\makecell[l]{DSI\_HOST\_GEN\_CMD\_WR\_ERR}} & 系统尝试通过通用接口写入命令，但 FIFO 已满。因此，命令未写入 & 重置 DSI 主机并重复传输 \\\hline
    \hyperref[fielddesc:DSIHOSTGENPLDWRERR]{\makecell[l]{DSI\_HOST\_GEN\_PLD\_WR\_ERR}} & 系统尝试通过通用接口写入有效载荷数据，但 FIFO 已满。因此，有效载荷未写入 & 重置 DSI 主机并重复传输 \\\hline
    \hyperref[fielddesc:DSIHOSTGENPLDSENDERR]{\makecell[l]{DSI\_HOST\_GEN\_PLD\_SEND\\\_ERR}} & 在通用接口数据包构建期间，有效载荷 FIFO 变空且发送了损坏的数据 & 重置 DSI 主机并重复传输 \\\hline
    \hyperref[fielddesc:DSIHOSTGENPLDRDERR]{\makecell[l]{DSI\_HOST\_GEN\_PLD\_RD\_ERR}} & 在 DCS 读取数据期间，有效载荷 FIFO 变空且发送到接口的数据已损坏 & 重置 DSI 主机并重复传输 \\\hline
    \hyperref[fielddesc:DSIHOSTGENPLDRECEVERR]{\makecell[l]{DSI\_HOST\_GEN\_PLD\_RECEV\\\_ERR}} & 在通用接口数据包读回期间，有效载荷 FIFO 变满且接收到的数据已损坏 & 重置 DSI 主机并重复传输 \\\hline
    \hyperref[fielddesc:DSIHOSTDPIBUFFPLDUNDER]{\makecell[l]{DSI\_HOST\_DPI\_BUFF\_PLD\\\_UNDER}} & 在读取有效载荷以构建视频模式的 DSI 数据包时发生了下溢 & 重置 DSI 主机并重复传输 \\\hline

\end{longtable}

%%%%%%%%%%%%%%%%%%%%%%%%%%%%%%
%%% 配置过程 %%%
%%%%%%%%%%%%%%%%%%%%%%%%%%%%%%

\section{配置过程}
\label{sec:mipidsi-prog-procedures}

\subsection{DSI 桥接配置}
\label{sec:dsi-bridge-configuration}
\begin{itemize}
    \item 配置 DSI Bridge 时钟并释放复位。
    \item 配置输入和输出颜色模式。详细信息请参阅 \hyperref[sec:dsi-bridge]{DSI Bridge}。
    \item 配置视频时序。有关详细信息，请参阅 \hyperref[sec:clocks-and-display-resolution]{时钟和显示分辨率}。
    \item 将 0 写入 \hyperref[fielddesc:DSIBRGDSIDMAFLOWCONTROLLER]{DSI\_BRG\_DSI\_DMA\_FLOW\_CONTROLLER} 设置 VDMA 为流控制器。
    \item 配置 \hyperref[fielddesc:DSIBRGDMABURSTLEN]{DSI\_BRG\_DMA\_BURST\_LEN} DMA 突发长度以匹配 VDMA 的配置。
    \item 配置 \hyperref[fielddesc:DSIBRGDSIRAWBUFALMOSTEMPTYTHRD]{DSI\_BRG\_DSI\_RAW\_BUF\_ALMOST\_EMPTY\_THRD} 为 $1024 - \hyperref[fielddesc:DSIBRGDMABURSTLEN]{DSI\_BRG\_DMA\_BURST\_LEN}$ 获得更好的性能。
    \item 将 1 写入 \hyperref[fielddesc:DSIBRGDPIEN]{DSI\_BRG\_DPI\_EN} 以启用 DPI。
    \item 将 1 写入 \hyperref[fielddesc:DSIBRGDSIEN]{DSI\_BRG\_DSI\_EN} 以启用 DSI Bridge。
    \item 配置 VDMA 并将 VDMA 设置为流控制器。详细信息请参阅 \ref{mod:vdma}。
\end{itemize}
\label{sec:mipidsi-dsi-bridge-programming-proc}

\subsection{D-PHY 配置}
\label{sec:mipidsi-dsi-dphy-programming-proc}
\begin{itemize}
    \item 将 \hyperref[fielddesc:DSIHOSTPHYSHUTDOWNZ]{DSI\_HOST\_PHY\_SHUTDOWNZ} 和 \hyperref[fielddesc:DSIHOSTPHYRSTZ]{DSI\_HOST\_PHY\_RSTZ} 设置为 0，将 \hyperref[fielddesc:DSIHOSTPHYTESTCLR]{DSI\_HOST\_PHY\_TESTCLR} 设置为 1。
    \item 应用适当的 DSI\_DPHY\_CFG\_CLK 和 DSI\_DPHY\_PLL\_REFCLK。
    \item 通过 \hyperref[sec:test-code-0x44]{测试代码 0x44} 配置 hsfreqrange。详细值请参见表 \ref{tab:mipidsi-hsfreq}。
    \item 配置 PLL 参数。详细信息参见 \hyperref[sec:pll-parameters]{PLL 参数}。
    \item 配置 \hyperref[fielddesc:DSIHOSTTOCLKDIVISION]{DSI\_HOST\_TO\_CLK\_DIVISION} 以应用适当的 TX\_ESC\_CLK 分频比。
    \item 将 \hyperref[fielddesc:DSIHOSTPHYSHUTDOWNZ]{DSI\_HOST\_PHY\_SHUTDOWNZ} 和 \hyperref[fielddesc:DSIHOSTPHYRSTZ]{DSI\_HOST\_PHY\_RSTZ} 设置为 1。
    \item 等待 \hyperref[fielddesc:DSIHOSTPHYSTOPSTATE0LANE]{DSI\_HOST\_PHY\_STOPSTATE0LANE} 和 \hyperref[fielddesc:DSIHOSTPHYSTOPSTATE1LANE]{DSI\_HOST\_PHY\_STOPSTATE1LANE} 信号置位。
\end{itemize}

\subsection{MIPI DSI 通用配置}
\label{sec:mipidsi-dsi-common-programming-proc}
\begin{itemize}
    \item 释放全局复位。
    \item 配置 \hyperref[fielddesc:DSIHOSTNLANES]{DSI\_HOST\_N\_LANES} 以定义要使用的通道数量。
    \item 配置视频时序设置。有关详细信息，请参阅 \hyperref[sec:clocks-and-display-resolution]{时钟和显示分辨率}。
    \item 通过 \hyperref[sec:mipidsi-dsi-dphy-programming-proc]{D-PHY 配置} 初始化 DPHY。
    \item 将 \hyperref[fielddesc:DSIHOSTSHUTDOWNZ]{DSI\_HOST\_SHUTDOWNZ} 写入 1 以释放 DSI 主机的关机状态。
\end{itemize}

\subsection{系统内存数据输出流配置}
\begin{itemize}
    \item 参考 \hyperref[sec:dsi-bridge-configuration]{DSI Bridge 配置} 中的方法配置 DSI Bridge。
    \item 参考 \hyperref[sec:mipidsi-dsi-common-programming-proc]{MIPI DSI 通用配置} 配置 DSI Host。
\end{itemize}

\subsection{视频模式输出流配置}
\begin{itemize}
    \item 通过寄存器字段 \hyperref[fielddesc:DSIHOSTVPGORIENTATION]{DSI\_HOST\_VPG\_ORIENTATION}、\hyperref[fielddesc:DSIHOSTVPGMODE]{DSI\_HOST\_VPG\_MODE} 和 \hyperref[fielddesc:DSIHOSTVPGEN]{DSI\_HOST\_VPG\_EN} 配置模式。有关详细信息，请参阅 \hyperref[sec:video-pattern-generator]{视频图案生成器}。
    \item 通过 \hyperref[sec:mipidsi-dsi-common-programming-proc]{MIPI DSI 通用配置} 配置 DSI 主机。
\end{itemize}


\subsection{使用 DCS 命令操作显示寄存器}
\begin{itemize}
    \item 写入显示
        \begin{itemize}
            \item 将数据包有效载荷写入 \hyperref[regdesc:DSIHOSTGENPLDDATAREG]{DSI\_HOST\_GEN\_PLD\_DATA\_REG}。
            \item 将数据包头写入 \hyperref[regdesc:DSIHOSTGENHDRREG]{DSI\_HOST\_GEN\_HDR\_REG}。
        \end{itemize}
    \item 从显示中读取
        \begin{itemize}
            \item 将数据包头写入 \hyperref[regdesc:DSIHOSTGENHDRREG]{DSI\_HOST\_GEN\_HDR\_REG}。
            \item 从 \hyperref[regdesc:DSIHOSTGENPLDDATAREG]{DSI\_HOST\_GEN\_PLD\_DATA\_REG} 读取响应数据。
        \end{itemize}
\end{itemize}


%%%%%%%%%%%%%%%%%%%%%%%%
%%% 寄存器摘要 %%%
%%%%%%%%%%%%%%%%%%%%%%%%


\newpage
\hypertarget{mipidsi-reg-summ}{}
\section{寄存器列表}
\label{sec:mipidsi-reg-summ}

\subsection{MIPI DSI 桥接寄存器列表}
本节中的地址为相对于 MIPI DSI BRIDGE 的基地址，该地址在第 \ref{mod:sysmem} 章的表 \ref{tab:sysmem-base-address} \textit{\nameref{mod:sysmem}} 中提供。

\textbf{Access} 列中的缩写在 \hyperref[glossary-access-types]{\textit{寄存器的访问类型}} 部分进行了说明。

\subfile{\modulefiles/42-MIPIDSI-BRG-reg-summ__CN}

\subsection{MIPI DSI 主机寄存器列表}
本节中的地址为相对于 MIPI DSI HOST 基地址，该地址在第 \ref{mod:sysmem} 章的表 \ref{tab:sysmem-base-address} 中提供 \textit{\nameref{mod:sysmem}}。

列 \textbf{访问}中给出的缩写在 \hyperref[glossary-access-types]{\textit{寄存器的访问类型}} 部分中进行了说明。

\subfile{\modulefiles/42-MIPIDSI-reg-summ__CN}



%%%%%%%%%%%%%%%%%
%%% 寄存器 %%%
%%%%%%%%%%%%%%%%%


\newpage

\section{寄存器}

\subsection{MIPI DSI 桥接寄存器}
本节中的地址相对于表 \ref{tab:sysmem-base-address} 中提供的 MIPI DSI Bridge 基地址，该表位于第 \ref{mod:sysmem} 章 \textit{\nameref{mod:sysmem}}。

有关如何配置保留字段，请参阅 \hyperref[programming-reserved-field]{配置保留寄存器字段} 部分。

\subfile{\modulefiles/42-MIPIDSI-BRG-reg__CN}

\subsection{MIPI DSI 主机寄存器}
本节中的地址相对于 MIPI DSI Host 基地址，该地址在第 \ref{mod:sysmem} 章的表 \ref{tab:sysmem-base-address} 中提供 \textit{\nameref{mod:sysmem}}。

有关如何配置保留字段，请参阅 \hyperref[programming-reserved-field]{配置保留寄存器字段} 部分。

\subfile{\modulefiles/42-MIPIDSI-reg__CN}

\end{document}
